% Candidate rules for offers, prices and conditions.  Written by hand; see docs/one-order-coverage.md.

\section{Offers, prices and conditions}

% Second survey round: candidates written from the survey notes in docs/one-order-survey.md.

\begin{theorem}[Offer expiry ends other limits (C-24)]
  \label{thm:offers-offer-expiry-ends-limits}
  \lean{Rules.OneOrder.Offers.guarantees_end_by_expiry}
  \uses{def:offer, def:instant, thm:accept-rejects-expired, def:offers-price-guarantee}
  Before an order exists, no price or inventory guarantee on an offer
  outlasts the offer's expiry. Once the offer's expiry instant has passed,
  both have lapsed too, unless an order was created from it. A payment time
  limit takes effect only once an order exists.

  Case to check: An offer item whose price guarantee ends an hour after the
  offer expires. A model that reads the guarantee field alone treats the
  price as guaranteed for an hour after the offer is gone.

  Open point: The Offers page also keeps the other limits alive when an
  order was changed, while the Time Limits page names only order creation.
  The business must decide whether the expiry of a reshop offer ends limits
  that would apply to the existing order it changes.

  Sources: \href{https://standards.atlassian.net/wiki/spaces/EASD/pages/600637464/Offers}{EASD guide, Offers};
  \href{https://standards.atlassian.net/wiki/spaces/EASD/pages/574423066/Time+Limits}{EASD guide, Time Limits}.
\end{theorem}
\begin{proof}
  \uses{def:offers-price-guarantee}
  A guarantee runs until the earlier of its own limit and the offer's
  expiry, so under either reading it holds at no instant after the expiry.
\end{proof}

\begin{definition}[Price guarantee at an instant (C-25)]
  \label{def:offers-price-guarantee}
  \lean{Rules.OneOrder.Offers.guaranteesAt, Rules.OneOrder.Offers.GuaranteeEnd, Rules.OneOrder.Offers.Guarantees, Rules.OneOrder.Offers.TimedOffer, Rules.OneOrder.Offers.TimedOfferItem, Rules.OneOrder.Offers.ItemLimits, Rules.OneOrder.Offers.TimedOffer.holds}
  \uses{def:offer-item, def:instant}
  An offer item's price is guaranteed at an instant only when the item
  carries a price guarantee time limit and the instant is inside both that
  limit and the offer's expiry. With no such limit the price is
  unguaranteed. The offer's expiry on its own guarantees neither price nor
  inventory. Whether a limit's own instant is inside is a reading of the
  guide, which the model takes as a parameter.

  Case to check: An item whose price guarantee ends at the same instant as
  the offer. Accepted at that instant, the offer is rejected as expired, yet
  an inclusive reading of the guarantee still calls the price guaranteed.
  The source does not fix the boundary.

  Sources: \href{https://standards.atlassian.net/wiki/spaces/EASD/pages/574423066/Time+Limits}{EASD guide, Time Limits};
  \href{https://standards.atlassian.net/wiki/spaces/EASD/pages/600637464/Offers}{EASD guide, Offers}.
\end{definition}

\begin{theorem}[No price guarantee without a limit]
  \label{thm:offers-no-limit-no-guarantee}
  \lean{Rules.OneOrder.Offers.no_limit_no_guarantee}
  \uses{def:offers-price-guarantee}
  Under every reading, an offer item with no price guarantee time limit has
  no guaranteed price at any instant.

  Source: \href{https://standards.atlassian.net/wiki/spaces/EASD/pages/574423066/Time+Limits}{EASD guide, Time Limits}.
\end{theorem}
\begin{proof}
  \uses{def:offers-price-guarantee}
  Unfold the rule: a missing limit holds nothing.
\end{proof}

\begin{theorem}[Price guaranteed before its limit]
  \label{thm:offers-price-guaranteed-before}
  \lean{Rules.OneOrder.Offers.price_guaranteed_before}
  \uses{def:offers-price-guarantee}
  Under every reading, an offer item's price is guaranteed at every instant
  before both its price guarantee time limit and the offer's expiry.

  Source: \href{https://standards.atlassian.net/wiki/spaces/EASD/pages/574423066/Time+Limits}{EASD guide, Time Limits}.
\end{theorem}
\begin{proof}
  \uses{def:offers-price-guarantee}
  An instant before both is before the earlier of the two.
\end{proof}

\begin{theorem}[Inventory guaranteed before its limit]
  \label{thm:offers-inventory-guaranteed-before}
  \lean{Rules.OneOrder.Offers.inventory_guaranteed_before}
  \uses{def:offers-price-guarantee}
  Under every reading, an offer item's inventory is guaranteed at every
  instant before both its inventory guarantee time limit and the offer's
  expiry.

  Source: \href{https://standards.atlassian.net/wiki/spaces/EASD/pages/574423066/Time+Limits}{EASD guide, Time Limits}.
\end{theorem}
\begin{proof}
  \uses{def:offers-price-guarantee}
  An instant before both is before the earlier of the two.
\end{proof}

\begin{theorem}[Guarantee ending at its instant]
  \label{thm:offers-guarantee-exclusive-acceptable}
  \lean{Rules.OneOrder.Offers.guaranteed_while_acceptable_of_exclusive}
  \uses{def:offers-price-guarantee, def:one-accept}
  Under the reading that a guarantee ends at its instant, an offer item's
  price is guaranteed only at instants at which the offer can still be
  accepted.

  Source: \href{https://standards.atlassian.net/wiki/spaces/EASD/pages/574423066/Time+Limits}{EASD guide, Time Limits}.
\end{theorem}
\begin{proof}
  \uses{def:offers-price-guarantee, def:one-accept}
  A guaranteed price means the instant is before the offer's expiry, and
  acceptance is rejected as expired only at or after it.
\end{proof}

\begin{theorem}[Guaranteed prices only on offers that can be accepted]
  \label{thm:offers-guaranteed-while-acceptable}
  \lean{Rules.OneOrder.Offers.guaranteed_while_acceptable}
  \uses{def:offers-price-guarantee, def:one-accept, thm:offers-guarantee-exclusive-acceptable}
  Under every reading, an offer item's price is guaranteed only at instants
  at which the offer can still be accepted.

  Counterexample: an offer valid until 12:00 whose flight's price guarantee
  also ends at 12:00. At 12:00 the offer is rejected as expired. Under the
  reading that a guarantee still holds at its instant, the price is
  guaranteed at that instant. The example \texttt{guaranteeAtExpiry}
  computes both.

  Open point: does a price guarantee still hold at the instant its time
  limit names?

  Source: \href{https://standards.atlassian.net/wiki/spaces/EASD/pages/574423066/Time+Limits}{EASD guide, Time Limits}.
\end{theorem}
\begin{proof}
  \uses{def:offers-price-guarantee}
  Not proven. The counterexample above shows the statement is false for the
  rule as written.
\end{proof}

\begin{definition}[Accepting after the guarantee lapses (C-26)]
  \label{def:offers-accept-after-guarantee-lapse}
  \lean{Rules.OneOrder.Offers.LapsedPrice, Rules.OneOrder.Offers.TimedOffer.guaranteeLapsed}
  \uses{def:accept, thm:accept-price-eq-sum, thm:accept-rejects-expired, def:offers-price-guarantee}
  When a customer accepts an unexpired offer after an item's price guarantee
  has lapsed, the airline decides the order price for that item. The source
  recommends repricing into a new offer before payment. The model takes the
  answer as a parameter: the acceptance is refused and sent back for a new
  offer, or the order is created at the offered price.

  Case to check: thm:accept-price-eq-sum makes the order price equal the
  offer item prices. An offer with a 30 minute price guarantee accepted two
  hours later, one minute before expiry, is sold at a price the airline no
  longer stands behind.

  Open point: The business must decide whether acceptance in this window is
  rejected, forces a reprice into a new offer, or creates the order at the
  offered price. If repricing at order creation is allowed, the order price
  can differ from the offer price.

  Sources: \href{https://standards.atlassian.net/wiki/spaces/EASD/pages/574423066/Time+Limits}{EASD guide, Time Limits};
  \href{https://standards.atlassian.net/wiki/spaces/EASD/pages/574423098/Pricing+Offers}{EASD guide, Pricing Offers}.
\end{definition}

\begin{theorem}[Order time limits only extend (C-27)]
  \label{thm:offers-order-limits-only-extend}
  \uses{def:accept, def:instant, def:c11-payment-time-limit, def:offers-price-guarantee}
  When an order is created, the time left on an offer item's price guarantee
  carries into the matching order item. After that the airline may move the
  order item's price guarantee or payment time limit later and never
  earlier.

  Case to check: An order change that recomputes the payment time limit from
  the change instant with a shorter duration than the original would move
  the limit earlier.

  Source: \href{https://standards.atlassian.net/wiki/spaces/EASD/pages/574423066/Time+Limits}{EASD guide, Time Limits}.
\end{theorem}

\begin{definition}[Payment time limit at creation (C-28)]
  \label{def:offers-payment-limit-at-creation}
  \lean{Rules.OneOrder.Offers.paymentLimitAt, Rules.OneOrder.Offers.BothForms, Rules.OneOrder.Offers.secondsAfter}
  \uses{def:accept, def:instant, def:c11-payment-time-limit, def:c12-zero-commitment}
  An offer item states its payment time limit as a duration or as a fixed
  date and time. At order creation the order item's limit becomes the
  creation instant plus the duration, or the fixed date and time. A zero
  duration means payment is due at creation. Which form applies when an
  offer item states both is a reading of the guide, which the model takes
  as a parameter.

  Case to check: A fixed offer limit that falls before the order is created
  gives an order item that is overdue the moment it exists.

  Open point: The source does not say which form wins if an offer item
  carries both, or how long an order item with no payment time limit may
  stay unpaid. The business must decide both.

  Source: \href{https://standards.atlassian.net/wiki/spaces/EASD/pages/574423066/Time+Limits}{EASD guide, Time Limits}.
\end{definition}

\begin{definition}[Time limits of a new order]
  \label{def:offers-create-order}
  \lean{Rules.OneOrder.Offers.createOrder, Rules.OneOrder.Offers.CreationReading, Rules.OneOrder.Offers.NewOrder, Rules.OneOrder.Offers.NewOrderItem, Rules.OneOrder.Offers.CreateRejection, Rules.OneOrder.Offers.newItem, Rules.OneOrder.Offers.NewOrder.toOrder, Rules.OneOrder.Offers.NewOrder.payable, Rules.OneOrder.Offers.TimedOffer.toOffer}
  \uses{def:one-accept, def:offers-price-guarantee, def:offers-accept-after-guarantee-lapse, def:offers-payment-limit-at-creation, def:c11-payment-time-limit}
  Accepting an offer whose items carry time limits runs the checks of
  ONE-001, then, under the reading that asks for a new offer, refuses the
  acceptance if an item's price guarantee has lapsed. Each order item takes
  its payment time limit and its price guarantee from its offer item.

  Source: \href{https://standards.atlassian.net/wiki/spaces/EASD/pages/574423066/Time+Limits}{EASD guide, Time Limits}.
\end{definition}

\begin{theorem}[New order matches ONE-001]
  \label{thm:offers-create-order-is-accept}
  \lean{Rules.OneOrder.Offers.createOrder_toOrder}
  \uses{def:offers-create-order, def:one-accept}
  An order created from a timed offer, without its time limits, is the order
  ONE-001 creates from the same offer, so every guarantee of ONE-001 holds
  for it.

  Source: \href{https://standards.atlassian.net/wiki/spaces/EASD/pages/574423066/Time+Limits}{EASD guide, Time Limits}.
\end{theorem}
\begin{proof}
  \uses{def:offers-create-order, def:one-accept}
  The rule creates the order only when ONE-001 does, with the same items.
\end{proof}

\begin{theorem}[Payment time limit from a duration]
  \label{thm:offers-create-duration}
  \lean{Rules.OneOrder.Offers.createOrder_duration}
  \uses{def:offers-create-order, def:offers-payment-limit-at-creation}
  An order item created from an offer item whose payment time limit is
  stated only as a duration must be paid by the creation instant plus that
  duration.

  Source: \href{https://standards.atlassian.net/wiki/spaces/EASD/pages/574423066/Time+Limits}{EASD guide, Time Limits}.
\end{theorem}
\begin{proof}
  \uses{def:offers-create-order}
  Unfold the rule for that offer item.
\end{proof}

\begin{theorem}[Payment time limit from a fixed instant]
  \label{thm:offers-create-fixed}
  \lean{Rules.OneOrder.Offers.createOrder_fixed}
  \uses{def:offers-create-order, def:offers-payment-limit-at-creation}
  An order item created from an offer item whose payment time limit is
  stated only as a fixed date and time must be paid by that instant.

  Source: \href{https://standards.atlassian.net/wiki/spaces/EASD/pages/574423066/Time+Limits}{EASD guide, Time Limits}.
\end{theorem}
\begin{proof}
  \uses{def:offers-create-order}
  Unfold the rule for that offer item.
\end{proof}

\begin{theorem}[Price guarantee carried into the order]
  \label{thm:offers-create-keeps-guarantee}
  \lean{Rules.OneOrder.Offers.createOrder_keeps_guarantee}
  \uses{def:offers-create-order}
  Each order item keeps the price guarantee time limit of its offer item.
  This is the first half of C-27; moving limits after creation is not
  modelled.

  Source: \href{https://standards.atlassian.net/wiki/spaces/EASD/pages/574423066/Time+Limits}{EASD guide, Time Limits}.
\end{theorem}
\begin{proof}
  \uses{def:offers-create-order}
  The order items are the offer items mapped one by one, each keeping its
  limit.
\end{proof}

\begin{theorem}[No order after a lapse when repricing]
  \label{thm:offers-create-reprice}
  \lean{Rules.OneOrder.Offers.createOrder_guaranteed_of_reprice}
  \uses{def:offers-create-order, def:offers-accept-after-guarantee-lapse}
  Under the reading that asks for a new offer, no order is created while
  the price guarantee of one of the offer's items has lapsed.

  Source: \href{https://standards.atlassian.net/wiki/spaces/EASD/pages/574423098/Pricing+Offers}{EASD guide, Pricing Offers}.
\end{theorem}
\begin{proof}
  \uses{def:offers-create-order}
  The rule refuses the acceptance when any item's guarantee has lapsed.
\end{proof}

\begin{theorem}[No order at a lapsed price]
  \label{thm:offers-create-price-guaranteed}
  \lean{Rules.OneOrder.Offers.createOrder_price_guaranteed}
  \uses{def:offers-create-order, thm:offers-create-reprice}
  Under every reading, no order is created at the offered price while the
  price guarantee of one of the offer's items has lapsed.

  Counterexample: an offer valid from 10:00 to 12:00 whose flight's price
  is guaranteed until 10:30, accepted at 11:59. Under the reading that
  keeps the offered price, the order is created with the flight at
  300.00 EUR. The example \texttt{lapsedGuaranteeAccepted} computes both
  readings.

  Open point: is acceptance after a lapsed price guarantee refused, sent
  back for a new offer, or honoured at the offered price?

  Sources: \href{https://standards.atlassian.net/wiki/spaces/EASD/pages/574423066/Time+Limits}{EASD guide, Time Limits};
  \href{https://standards.atlassian.net/wiki/spaces/EASD/pages/574423098/Pricing+Offers}{EASD guide, Pricing Offers}.
\end{theorem}
\begin{proof}
  \uses{def:offers-create-order}
  Not proven. The counterexample above shows the statement is false for the
  rule as written.
\end{proof}

\begin{theorem}[No order item overdue at creation]
  \label{thm:offers-create-limit-after-creation}
  \lean{Rules.OneOrder.Offers.createOrder_limit_not_before_creation}
  \uses{def:offers-create-order, def:offers-payment-limit-at-creation}
  No order item starts with a payment time limit before the order's
  creation.

  Counterexample: an offer valid until 12:00 whose bag must be paid by
  10:30, accepted at 11:00. The bag's order item must be paid half an hour
  before the order existed. The example \texttt{fixedLimitBeforeCreation}
  computes it.

  Open point: is an acceptance after a fixed payment time limit refused, or
  does the order item get a new payment time limit?

  Source: \href{https://standards.atlassian.net/wiki/spaces/EASD/pages/574423066/Time+Limits}{EASD guide, Time Limits}.
\end{theorem}
\begin{proof}
  \uses{def:offers-create-order}
  Not proven. The counterexample above shows the statement is false for the
  rule as written.
\end{proof}

\begin{theorem}[Payment due at creation can be made]
  \label{thm:offers-zero-duration-payable}
  \lean{Rules.OneOrder.Offers.zero_duration_payable}
  \uses{def:offers-create-order, def:offers-payment-limit-at-creation, def:payment-accept}
  An order item whose offer item's payment time limit is a duration of zero
  can be paid at the instant the order is created.

  Counterexample: a flight at 300.00 EUR with a payment time limit of zero,
  accepted at 11:00 and paid in full by card at 11:00. The order item's
  limit is 11:00, and the payment rule of ONE-304 rejects a payment made at
  or after the limit. The example \texttt{zeroDurationPayment} computes it.

  Open point: is a payment made at the creation instant on time, or must it
  come with the request that creates the order?

  Source: \href{https://standards.atlassian.net/wiki/spaces/EASD/pages/574423066/Time+Limits}{EASD guide, Time Limits}.
\end{theorem}
\begin{proof}
  \uses{def:offers-create-order, def:payment-accept}
  Not proven. The counterexample above shows the statement is false for the
  rule as written.
\end{proof}

\begin{theorem}[Nesting of guarantee time limits (C-29)]
  \label{thm:offers-guarantee-limits-nest}
  \uses{def:offers-price-guarantee, def:c11-payment-time-limit}
  \notready
  When an item has a price guarantee time limit, that limit is at or after
  every separate guarantee limit on its base amount, each tax and each fee.
  Those part limits count only when the overall limit is present. The price
  guarantee should also end no later than the payment time limit.

  Case to check: The source's table guarantees the base amount and one tax
  for 5 days in the order. If the order item's overall guarantee stays at
  the 30 minutes carried from the offer, the parts outlive the whole.

  Open point: The source calls the overall limit the maximum of the parts,
  which reads either as equal to the latest part or as an upper bound, and
  the payment limit bound is only a recommendation. The business must pick
  one reading and decide whether a price guarantee past the payment limit is
  cut back or refused.

  Sources: \href{https://standards.atlassian.net/wiki/spaces/EASD/pages/574423066/Time+Limits}{EASD guide, Time Limits};
  \href{https://standards.atlassian.net/wiki/spaces/EASD/pages/574914636/Fees}{EASD guide, Fees}.
\end{theorem}

\begin{definition}[Restriction rule in force (C-30)]
  \label{def:offers-restriction-in-force}
  \lean{Rules.OneOrder.Offers.ruleInForce, Rules.OneOrder.Offers.Restriction, Rules.OneOrder.Offers.Stage, Rules.OneOrder.Offers.Journey, Rules.OneOrder.Offers.Journey.stageAt, Rules.OneOrder.Offers.Action, Rules.OneOrder.Offers.ChangeKind, Rules.OneOrder.Offers.ExpiryReading, Rules.OneOrder.Offers.Restriction.live, Rules.OneOrder.Offers.latestOf, Rules.OneOrder.Offers.candidates}
  \uses{def:offer-item, def:instant}
  Each change or cancel rule on an offer item belongs to one journey stage
  (before payment, before departure, after departure, no show) and starts at
  its effective date. At a given instant the rule in force is the one in the
  current stage with the latest effective date not after that instant. A
  rule stops when its stage ends or a later rule takes over. Whether a
  rule's expiration date also ends it is a reading of the guide, which the
  model takes as a parameter.

  Case to check: In the source example (payment 1 Jan 2019, departure 1 Feb
  2019) the no-change rule described as 7 days before departure starts on 24
  Jan, which is 8 days before. The text gives the middle fee as 1 CHF where
  the table and diagram give 60 CHF. Another example sets equal effective
  and expiration dates, so that rule is never in force.

  Open point: The message examples also give rules an expiration date, which
  the stage and override reading never uses. The business must decide which
  governs, and what applies before the first effective date or in a gap
  between one rule's expiration and the next rule. The source also does not
  say which of two rules of a stage applies when they start at the same
  instant.

  Sources: \href{https://standards.atlassian.net/wiki/spaces/EASD/pages/574259692/Offer+Restrictions}{EASD guide, Offer Restrictions};
  \href{https://standards.atlassian.net/wiki/spaces/EASD/pages/574816307/Offer+Conditions+and+Restrictions}{EASD guide, Offer Conditions and Restrictions}.
\end{definition}

\begin{theorem}[Rule in force belongs to the current stage]
  \label{thm:offers-rule-stage}
  \lean{Rules.OneOrder.Offers.ruleInForce_stage}
  \uses{def:offers-restriction-in-force}
  The rule in force for an action at an instant is one of the rules for
  that action and for the stage the journey is in at that instant.

  Source: \href{https://standards.atlassian.net/wiki/spaces/EASD/pages/574259692/Offer+Restrictions}{EASD guide, Offer Restrictions}.
\end{theorem}
\begin{proof}
  \uses{def:offers-restriction-in-force}
  The rule picks from the rules filtered by action and stage.
\end{proof}

\begin{theorem}[Rule in force has started]
  \label{thm:offers-rule-started}
  \lean{Rules.OneOrder.Offers.ruleInForce_started}
  \uses{def:offers-restriction-in-force}
  The rule in force at an instant has its effective date at or before that
  instant.

  Source: \href{https://standards.atlassian.net/wiki/spaces/EASD/pages/574259692/Offer+Restrictions}{EASD guide, Offer Restrictions}.
\end{theorem}
\begin{proof}
  \uses{def:offers-restriction-in-force}
  The rule picks only from rules that have started.
\end{proof}

\begin{theorem}[Latest started rule is in force]
  \label{thm:offers-rule-latest}
  \lean{Rules.OneOrder.Offers.ruleInForce_latest}
  \uses{def:offers-restriction-in-force}
  When rules for an action and the current stage have started, and have not
  expired where expiration dates end rules, one of them is in force, and
  none of them starts later than it.

  Source: \href{https://standards.atlassian.net/wiki/spaces/EASD/pages/574259692/Offer+Restrictions}{EASD guide, Offer Restrictions}.
\end{theorem}
\begin{proof}
  \uses{def:offers-restriction-in-force}
  By induction on the list of rules: the latest of a list is one of them and
  starts no earlier than any of them.
\end{proof}

\begin{theorem}[Expired rules are not in force]
  \label{thm:offers-rule-not-expired}
  \lean{Rules.OneOrder.Offers.ruleInForce_not_expired}
  \uses{def:offers-restriction-in-force}
  Under the reading that expiration dates end rules, the rule in force at an
  instant has not reached its expiration date.

  Source: \href{https://standards.atlassian.net/wiki/spaces/EASD/pages/574816307/Offer+Conditions+and+Restrictions}{EASD guide, Offer Conditions and Restrictions}.
\end{theorem}
\begin{proof}
  \uses{def:offers-restriction-in-force}
  Under that reading the rule picks only from rules before their expiration
  date.
\end{proof}

\begin{theorem}[Every stage with rules has a rule in force]
  \label{thm:offers-rule-covers-stage}
  \lean{Rules.OneOrder.Offers.ruleInForce_covers_stage}
  \uses{def:offers-restriction-in-force, thm:offers-rule-latest}
  While a journey is in a stage for which an action has rules, a rule is in
  force for that action.

  Counterexample: a journey paid at 10:00 on 1 January whose only flight
  change rule before departure starts on 2 January has no rule in force at
  18:00 on 1 January, under either reading. Under the reading that
  expiration dates end rules, a rule whose effective and expiration dates
  are equal is never in force. The example \texttt{noRuleInForce} computes
  both.

  Open point: what applies to a change or cancellation when no rule of the
  current stage is in force?

  Sources: \href{https://standards.atlassian.net/wiki/spaces/EASD/pages/574259692/Offer+Restrictions}{EASD guide, Offer Restrictions};
  \href{https://standards.atlassian.net/wiki/spaces/EASD/pages/574816307/Offer+Conditions+and+Restrictions}{EASD guide, Offer Conditions and Restrictions}.
\end{theorem}
\begin{proof}
  \uses{def:offers-restriction-in-force}
  Not proven. The counterexample above shows the statement is false for the
  rule as written.
\end{proof}

\begin{definition}[Voluntary change under the rule in force]
  \label{def:offers-voluntary-change}
  \lean{Rules.OneOrder.Offers.voluntaryChange, Rules.OneOrder.Offers.ChangeRequest, Rules.OneOrder.Offers.ChangeDecision, Rules.OneOrder.Offers.Restriction.feeOr}
  \uses{def:offers-restriction-in-force, def:change-cancellation-price-differential}
  A voluntary change is refused when the rule in force forbids it. When the
  rule allows it, the change goes to the airline's pricing of ONE-201 with
  the old and new prices, the rule's fee as the penalty and nothing
  forfeited. When no rule is in force, the decision says so.

  Source: \href{https://standards.atlassian.net/wiki/spaces/EASD/pages/574259692/Offer+Restrictions}{EASD guide, Offer Restrictions}.
\end{definition}

\begin{theorem}[Forbidden changes are refused]
  \label{thm:offers-change-refused}
  \lean{Rules.OneOrder.Offers.change_refused}
  \uses{def:offers-voluntary-change}
  A change the rule in force forbids is refused.

  Source: \href{https://standards.atlassian.net/wiki/spaces/EASD/pages/574259692/Offer+Restrictions}{EASD guide, Offer Restrictions}.
\end{theorem}
\begin{proof}
  \uses{def:offers-voluntary-change}
  Unfold the rule.
\end{proof}

\begin{theorem}[Rule's fee is the penalty]
  \label{thm:offers-change-fee}
  \lean{Rules.OneOrder.Offers.change_fee_is_penalty}
  \uses{def:offers-voluntary-change}
  A change the rule in force allows with a fee is quoted with that fee as
  its penalty.

  Source: \href{https://standards.atlassian.net/wiki/spaces/EASD/pages/574259692/Offer+Restrictions}{EASD guide, Offer Restrictions}.
\end{theorem}
\begin{proof}
  \uses{def:offers-voluntary-change}
  Unfold the rule.
\end{proof}

\begin{theorem}[No fee, no penalty]
  \label{thm:offers-change-no-fee}
  \lean{Rules.OneOrder.Offers.change_no_fee}
  \uses{def:offers-voluntary-change}
  A change the rule in force allows without a fee is quoted with no
  penalty.

  Source: \href{https://standards.atlassian.net/wiki/spaces/EASD/pages/574259692/Offer+Restrictions}{EASD guide, Offer Restrictions}.
\end{theorem}
\begin{proof}
  \uses{def:offers-voluntary-change}
  Unfold the rule.
\end{proof}

\begin{theorem}[Voluntary change under a restriction (C-31)]
  \label{thm:offers-change-under-restriction}
  \lean{Rules.OneOrder.Offers.change_net_is_fee_plus_difference}
  \uses{def:offers-restriction-in-force, def:offers-voluntary-change, def:c18-change, def:change-cancellation-price-differential, def:change-cancellation-due-amounts}
  A voluntary change of a given type, such as flight or name, is rejected
  when the rule in force forbids it. When the rule allows it, the customer
  pays the rule's fee, or nothing if the rule states no fee, plus any fare
  difference. The model proves the payment for a fee no larger than the
  price of the items changed, when the seller does not ask to keep the value
  for later use.

  Case to check: The source example described as a no-show change with total
  loss of fare carries a name change rule with a 1 CHF fee and no stage, so
  its description and its data disagree on what the customer owes.

  Open point: The source does not say what applies when no rule covers the
  current stage, or whether a fee marked approximate binds when the change
  is made. The business must choose a default for an uncovered stage and
  decide how an approximate fee is settled.

  Sources: \href{https://standards.atlassian.net/wiki/spaces/EASD/pages/574259692/Offer+Restrictions}{EASD guide, Offer Restrictions};
  \href{https://standards.atlassian.net/wiki/spaces/EASD/pages/574816307/Offer+Conditions+and+Restrictions}{EASD guide, Offer Conditions and Restrictions}.
\end{theorem}
\begin{proof}
  \uses{def:offers-voluntary-change, def:change-cancellation-due-amounts}
  The quote's price difference is the new price less the old, and a fee no
  larger than the old price is charged in full under both readings of
  ONE-202, netted or separate.
\end{proof}

\begin{definition}[Refund owed on voluntary cancellation (C-32)]
  \label{def:offers-refund-on-cancellation}
  \lean{Rules.OneOrder.Offers.cancellationRefund, Rules.OneOrder.Offers.RefundTerms, Rules.OneOrder.Offers.RefundTerms.amount, Rules.OneOrder.Offers.FullOrFixed, Rules.OneOrder.Offers.RefundDecision, Rules.OneOrder.Offers.CancelledItem, Rules.OneOrder.Offers.CancelledItem.refundableCharges}
  \uses{def:offers-restriction-in-force, def:c15-cancel-order, thm:c16-penalty-net, def:refund}
  The cancel rule in force decides what a voluntary cancellation returns. A
  rule that forbids cancellation returns nothing. A rule that allows it
  returns a full or a fixed refund, paid to the original form of payment
  when the rule says so. Which amount a refund marked full and also given a
  fixed amount returns is a reading of the guide, which the model takes as
  a parameter.

  Case to check: Reading that example as the full item price and reading it
  as 1.00 CHF give two different refunds for the same rule.

  Open point: One source example marks a refund as full and also gives a
  fixed amount of 1.00 CHF. The business must decide which governs when both
  are present, and how the rule's refund relates to the penalty netting in
  thm:c16-penalty-net.

  Source: \href{https://standards.atlassian.net/wiki/spaces/EASD/pages/574259692/Offer+Restrictions}{EASD guide, Offer Restrictions}.
\end{definition}

\begin{theorem}[Forbidden cancellation returns nothing]
  \label{thm:offers-cancel-forbidden}
  \lean{Rules.OneOrder.Offers.cancel_forbidden_nothing}
  \uses{def:offers-refund-on-cancellation}
  A cancellation the rule in force forbids returns nothing.

  Source: \href{https://standards.atlassian.net/wiki/spaces/EASD/pages/574259692/Offer+Restrictions}{EASD guide, Offer Restrictions}.
\end{theorem}
\begin{proof}
  \uses{def:offers-refund-on-cancellation}
  Unfold the rule.
\end{proof}

\begin{theorem}[Full refund returns the whole price]
  \label{thm:offers-cancel-full}
  \lean{Rules.OneOrder.Offers.cancel_full_refund}
  \uses{def:offers-refund-on-cancellation}
  A cancellation the rule in force allows with a full refund and no fixed
  amount returns the item's whole price.

  Source: \href{https://standards.atlassian.net/wiki/spaces/EASD/pages/574259692/Offer+Restrictions}{EASD guide, Offer Restrictions}.
\end{theorem}
\begin{proof}
  \uses{def:offers-refund-on-cancellation}
  Unfold the rule.
\end{proof}

\begin{theorem}[Fixed refund returns its amount]
  \label{thm:offers-cancel-fixed}
  \lean{Rules.OneOrder.Offers.cancel_fixed_refund}
  \uses{def:offers-refund-on-cancellation}
  A cancellation the rule in force allows with a fixed refund not marked
  full returns that amount.

  Source: \href{https://standards.atlassian.net/wiki/spaces/EASD/pages/574259692/Offer+Restrictions}{EASD guide, Offer Restrictions}.
\end{theorem}
\begin{proof}
  \uses{def:offers-refund-on-cancellation}
  Unfold the rule.
\end{proof}

\begin{theorem}[Refund to the original form of payment]
  \label{thm:offers-cancel-original-form}
  \lean{Rules.OneOrder.Offers.cancel_to_original_form}
  \uses{def:offers-refund-on-cancellation}
  A cancellation the rule in force allows, with the refund sent to the
  original form of payment, is refunded to the original form of payment.

  Source: \href{https://standards.atlassian.net/wiki/spaces/EASD/pages/574259692/Offer+Restrictions}{EASD guide, Offer Restrictions}.
\end{theorem}
\begin{proof}
  \uses{def:offers-refund-on-cancellation}
  Unfold the rule.
\end{proof}

\begin{theorem}[Taxes on a non-refundable fare (C-33)]
  \label{thm:offers-refundable-taxes}
  \lean{Rules.OneOrder.Offers.cancel_returns_refundable_charges}
  \uses{def:offers-refund-on-cancellation, thm:offers-tax-summary-total, def:refund}
  Each tax, fee and charge on a price carries its own refund flag. When an
  unused item is cancelled under a rule that forfeits the fare, the customer
  is still owed the taxes, fees and charges flagged refundable.

  Counterexample: a trip at 500.00 EUR with a refundable airport tax of
  25.00 EUR, cancelled on 25 January under a rule that forbids
  cancellation, returns nothing; so does a cancellation after a no show.
  Under the reading that a refund marked full with a fixed amount of
  1.00 EUR returns the fixed amount, the cancellation returns 1.00 EUR. The
  examples \texttt{nonRefundableFareTaxes} and \texttt{fullMarkedWithAmount}
  compute them.

  Open point: The source gives refund flags per tax and per fee and refund
  rules per offer item, but never says how they combine. The business must
  decide whether a no-refund cancel rule also keeps refundable taxes, and
  whether a no show changes the answer.

  Sources: \href{https://standards.atlassian.net/wiki/spaces/EASD/pages/574455915/Taxes+Fees+and+Charges}{EASD guide, Taxes, Fees, and Charges};
  \href{https://standards.atlassian.net/wiki/spaces/EASD/pages/574914636/Fees}{EASD guide, Fees};
  \href{https://standards.atlassian.net/wiki/spaces/EASD/pages/574259692/Offer+Restrictions}{EASD guide, Offer Restrictions}.
\end{theorem}
\begin{proof}
  \uses{def:offers-refund-on-cancellation}
  Not proven. The counterexample above shows the statement is false for the
  rule as written.
\end{proof}

\begin{definition}[Totalling a price from its parts]
  \label{def:offers-price-parts}
  \lean{Rules.OneOrder.Offers.priceTotal, Rules.OneOrder.Offers.PriceParts, Rules.OneOrder.Offers.PriceTotal, Rules.OneOrder.Offers.PriceTotal.of, Rules.OneOrder.Offers.PriceRejection, Rules.OneOrder.Offers.Charge, Rules.OneOrder.Offers.ConsumptionTax, Rules.OneOrder.Offers.ExcludedTax, Rules.OneOrder.Offers.ExcludedTax.amount, Rules.OneOrder.Offers.PriceParts.lineExcluded, Rules.OneOrder.Offers.PriceParts.amounts, Rules.OneOrder.Offers.sumOf}
  \uses{def:money, def:money-add, def:offer-item}
  A price stated with its parts is rejected when its amounts are not all in
  one currency or its discount is larger than the amount before discounts.
  Otherwise its base amount is the amount before discounts less the
  discount, its tax summary total is the sum of its entries, and its total
  adds the base amount, the surcharges, the fees, the tax summary total and
  the consumption tax not included in it. Whether that tax comes from the
  lines marked as not included or from the stated total of excluded tax is
  a reading of the guide, which the model takes as a parameter. Discounts
  are stated as amounts.

  Sources: \href{https://standards.atlassian.net/wiki/spaces/EASD/pages/574455915/Taxes+Fees+and+Charges}{EASD guide, Taxes, Fees, and Charges};
  \href{https://standards.atlassian.net/wiki/spaces/EASD/pages/2588377182/Consumption+Taxes}{EASD guide, Consumption Taxes}.
\end{definition}

\begin{theorem}[Mixed-currency prices are rejected]
  \label{thm:offers-price-mixed-currency}
  \lean{Rules.OneOrder.Offers.price_rejects_mixed_currency}
  \uses{def:offers-price-parts}
  A price with an amount in a currency other than that of its amount before
  discounts is rejected.

  Source: \href{https://standards.atlassian.net/wiki/spaces/EASD/pages/749305864/Currency+in+Monetary+Amounts}{EASD guide, Currency in Monetary Amounts}.
\end{theorem}
\begin{proof}
  \uses{def:offers-price-parts}
  The currency check runs first.
\end{proof}

\begin{theorem}[Discount above the base is rejected]
  \label{thm:offers-price-large-discount}
  \lean{Rules.OneOrder.Offers.price_rejects_large_discount}
  \uses{def:offers-price-parts}
  A price whose discount is larger than its amount before discounts is never
  totalled.

  Source: \href{https://standards.atlassian.net/wiki/spaces/EASD/pages/574849038/Discounts}{EASD guide, Discounts}.
\end{theorem}
\begin{proof}
  \uses{def:offers-price-parts}
  Every branch the rule can reach for such a price is a rejection.
\end{proof}

\begin{theorem}[Base amount after discounts]
  \label{thm:offers-price-base}
  \lean{Rules.OneOrder.Offers.price_base}
  \uses{def:offers-price-parts}
  A totalled price's base amount is its amount before discounts less its
  discount.

  Source: \href{https://standards.atlassian.net/wiki/spaces/EASD/pages/574849038/Discounts}{EASD guide, Discounts}.
\end{theorem}
\begin{proof}
  \uses{def:offers-price-parts}
  Unfold the rule.
\end{proof}

\begin{theorem}[Tax summary total from its entries]
  \label{thm:offers-price-tax-total}
  \lean{Rules.OneOrder.Offers.price_tax_total}
  \uses{def:offers-price-parts}
  A totalled price's tax summary total is the sum of its entries.

  Source: \href{https://standards.atlassian.net/wiki/spaces/EASD/pages/574455915/Taxes+Fees+and+Charges}{EASD guide, Taxes, Fees, and Charges}.
\end{theorem}
\begin{proof}
  \uses{def:offers-price-parts}
  Unfold the rule.
\end{proof}

\begin{theorem}[Total from base, surcharges, fees and taxes]
  \label{thm:offers-price-total-adds}
  \lean{Rules.OneOrder.Offers.price_total_from_parts}
  \uses{def:offers-price-parts}
  A totalled price's total is its base amount plus its surcharges, its
  fees, its tax summary total and the consumption tax not included in it,
  as the reading takes that tax, all in the price's currency.

  Source: \href{https://standards.atlassian.net/wiki/spaces/EASD/pages/574455915/Taxes+Fees+and+Charges}{EASD guide, Taxes, Fees, and Charges}.
\end{theorem}
\begin{proof}
  \uses{def:offers-price-parts}
  Unfold the rule.
\end{proof}

\begin{theorem}[Excluded tax from the lines]
  \label{thm:offers-price-lines-reading}
  \lean{Rules.OneOrder.Offers.price_adds_excluded_lines_of_lines}
  \uses{def:offers-price-parts}
  Under the reading that takes the excluded consumption tax from the lines,
  a totalled price adds the tax of every line marked as not included.

  Source: \href{https://standards.atlassian.net/wiki/spaces/EASD/pages/2588377182/Consumption+Taxes}{EASD guide, Consumption Taxes}.
\end{theorem}
\begin{proof}
  \uses{def:offers-price-parts}
  Unfold the rule under that reading.
\end{proof}

\begin{theorem}[Price total from its parts (C-34)]
  \label{thm:offers-price-total-from-parts}
  \lean{Rules.OneOrder.Offers.price_adds_excluded_lines}
  \uses{def:money-add, def:offer-item, def:offers-price-parts, thm:offers-price-total-adds, thm:offers-price-lines-reading, thm:offers-tax-summary-total, thm:offers-consumption-tax-lines}
  A price's total equals its base amount plus surcharges, fees, the tax
  summary total and any consumption tax not already included. The base
  amount is stated after discounts: the pre-discount amount minus the
  discount, where a percentage discount is that share of the pre-discount
  amount.

  Case to check: In the Canadian example every HST line is marked as not
  included, yet each reports zero excluded tax. Adding the excluded tax
  raises the total by 58.78 CAD; trusting the zero totals leaves it
  unchanged.

  Counterexample: a base of 452.15 with an HST line of 58.78 marked as not
  included and a stated total of excluded tax of zero, in US dollars since
  the model has no Canadian dollar. Taking the tax from the lines gives a
  total of 510.93; trusting the stated total gives 452.15. The example
  \texttt{excludedTaxReadings} computes both.

  Open point: The source names the parts and gives one example that adds up
  (95.00 base and 20.00 tax give 115.00), but states no formula. The
  business must decide whether fees and surcharges sit inside the tax
  summary or add separately, whether a discount applies before taxes are
  computed, and whether excluded consumption tax adds to the total. The
  model states discounts as amounts, so a percentage discount and its
  rounding are not modelled.

  Sources: \href{https://standards.atlassian.net/wiki/spaces/EASD/pages/574455915/Taxes+Fees+and+Charges}{EASD guide, Taxes, Fees, and Charges};
  \href{https://standards.atlassian.net/wiki/spaces/EASD/pages/574914636/Fees}{EASD guide, Fees};
  \href{https://standards.atlassian.net/wiki/spaces/EASD/pages/574849038/Discounts}{EASD guide, Discounts};
  \href{https://standards.atlassian.net/wiki/spaces/EASD/pages/2588377182/Consumption+Taxes}{EASD guide, Consumption Taxes}.
\end{theorem}
\begin{proof}
  \uses{def:offers-price-parts}
  Not proven. The counterexample above shows the statement is false for the
  rule as written.
\end{proof}

\begin{theorem}[Tax summary total and detail (C-35)]
  \label{thm:offers-tax-summary-total}
  \uses{def:money-add, def:accept, def:order-item}
  A tax summary's total equals the sum of its individual tax, fee and charge
  amounts, all in one currency. A shopping response may give only the total,
  but the order created from an accepted offer must hold the individual
  entries, including consumption tax detail.

  Case to check: An offer shopped with only a tax total and accepted without
  a pricing step has no entries to copy into the order. The source does not
  say where they come from or require them to add up to the shopped total.

  Sources: \href{https://standards.atlassian.net/wiki/spaces/EASD/pages/574455915/Taxes+Fees+and+Charges}{EASD guide, Taxes, Fees, and Charges};
  \href{https://standards.atlassian.net/wiki/spaces/EASD/pages/2588377182/Consumption+Taxes}{EASD guide, Consumption Taxes}.
\end{theorem}

\begin{theorem}[Consumption tax lines add up (C-36)]
  \label{thm:offers-consumption-tax-lines}
  \uses{def:money, def:money-add, thm:offers-tax-summary-total}
  \notready
  In each consumption tax entry, the tax is the rate applied to the taxable
  base, and the entry's total tax is the sum of its line taxes. A gross
  amount equals its net amount plus every consumption tax charged on it. All
  amounts in an entry are in the price's currency.

  Case to check: Italian example: line tax 1.82 against entry total 1.80.
  Canadian example: line tax 0.93 against entry total 4.55. Australian
  example: gross in AUD with tax in EUR. Tahiti example: two taxes on one
  base, so a per-line check that gross equals net plus that line's tax
  fails.

  Open point: The source shows these relations only through examples,
  several of which break them, and gives no rounding rule. The business must
  decide which figure is authoritative when a line and its total disagree,
  and how to round.

  Source: \href{https://standards.atlassian.net/wiki/spaces/EASD/pages/2588377182/Consumption+Taxes}{EASD guide, Consumption Taxes}.
\end{theorem}

\begin{theorem}[Seller margin stays outside order (C-37)]
  \label{thm:offers-margin-outside-order}
  \uses{def:accept, thm:accept-price-eq-sum, def:payment}
  When the seller is merchant of record and adds a margin, the airline's
  order total equals the airline offer price including taxes, fees and
  charges. The order holds no record of the margin, and the seller pays the
  airline the offer price.

  Case to check: The source places the airline as merchant of record out of
  scope. In that case the payment recorded on the order would include the
  margin and exceed the order total, which a rule tying payments to the
  order total would reject.

  Source: \href{https://standards.atlassian.net/wiki/spaces/EASD/pages/574849050/Commission+and+Margin}{EASD guide, Commission and Margin}.
\end{theorem}

\begin{definition}[Mandatory and optional offer items (C-38)]
  \label{def:offers-mandatory-optional-items}
  \uses{def:offer, def:offer-item, def:accept, thm:c02-item-per-offer-item, thm:accept-price-eq-sum, thm:c04-item-unsplit}
  \notready
  Each offer item is mandatory or optional, and its services and price are
  fixed once offered. Accepting an offer takes every mandatory item whole
  plus the optional items chosen; a-la-carte items are all optional, priced
  per unit, and bought in a chosen quantity. Every offer other than an
  a-la-carte offer should hold at least one mandatory item.

  Case to check: thm:c02-item-per-offer-item creates one order item per
  offer item, which fails when the customer skips an optional item. The
  source's own figure shows an offer total of 80 USD equal to one item while
  a second item is listed.

  Open point: The source does not say whether an offer's total price counts
  its optional items. The business must decide that, and how acceptance
  records the chosen items and quantities.

  Sources: \href{https://standards.atlassian.net/wiki/spaces/EASD/pages/600637464/Offers}{EASD guide, Offers};
  \href{https://standards.atlassian.net/wiki/spaces/SRSIA/pages/574619905/General+Principles+and+Assumptions}{Interline guide, General Principles and Assumptions}.
\end{definition}

\begin{theorem}[Offer responsible airline owns order (C-39)]
  \label{thm:offers-ora-owns-order}
  \uses{def:offer, def:accept, def:order}
  Each offer has one Offer Responsible Airline, the airline that returns the
  combined offer; every other airline in it is a Participating Offer
  Airline. The order created from the offer is owned by the Offer
  Responsible Airline, which receives the seller's payment commitment and
  holds the single source of the order's data.

  Case to check: In the interline ordering flow each participating airline
  also keeps its own order for its services. Status read from that order can
  disagree with the owner's order unless every status change reaches the
  owner.

  Sources: \href{https://www.iata.org/contentassets/72cbd60393ff42b5975d90ce9e049a7d/oneorder-resolution-797.pdf}{IATA Resolution 797};
  \href{https://standards.atlassian.net/wiki/spaces/SWO/pages/574816490/Roles}{Settlement with Orders, Roles};
  \href{https://standards.atlassian.net/wiki/spaces/SRSIA/pages/574685353/Ordering}{Interline guide, Ordering}.
\end{theorem}

\begin{definition}[Terms of a combined interline offer (C-40)]
  \label{def:offers-interline-offer-terms}
  \uses{thm:offers-ora-owns-order, def:offers-restriction-in-force, def:offers-mandatory-optional-items}
  \notready
  When the Offer Responsible Airline combines a participating airline's
  offer item into its own offer, it keeps that item whole and accepts the
  item's settlement value, settlement terms and sales conditions. It sets
  the customer price itself, and the settlement value, held per service,
  stays between the two airlines.

  Case to check: A participating airline's segment is non-refundable while
  the combined offer item is shown as fully refundable. The customer gets a
  full refund and the Offer Responsible Airline still owes the settlement
  value.

  Open point: The source does not say how the participating airline's change
  and cancel rules become the restrictions shown to the customer, or whether
  the customer price must cover the settlement values. The business must
  decide both.

  Sources: \href{https://standards.atlassian.net/wiki/spaces/SRSIA/pages/574816454/Shopping}{Interline guide, Shopping};
  \href{https://standards.atlassian.net/wiki/spaces/SRSIA/pages/574619905/General+Principles+and+Assumptions}{Interline guide, General Principles and Assumptions};
  \href{https://standards.atlassian.net/wiki/spaces/SRSIA/pages/574914728/Glossary+of+Terms}{Interline guide, Glossary of Terms};
  \href{https://standards.atlassian.net/wiki/spaces/SRSIA/overview}{Interline guide}.
\end{definition}

\begin{definition}[Currency of an offer's amounts (C-41)]
  \label{def:offers-offer-currency}
  \uses{def:money, def:currency, def:offer, def:money-add}
  \notready
  Every amount carries a currency, zero included. Amounts that add into an
  offer's totals are in the offer's currency and use that currency's number
  of decimals (two for GBP, three for TND, none for JPY). Local or
  equivalent amounts in another currency stay out of totals.

  Case to check: The source's Tahiti example states XPF amounts to two
  decimals, although ISO 4217 gives XPF none. A minor-unit model reads
  100000.00 XPF as either 100000 or 10000000 units depending on which rule
  it trusts.

  Open point: The source lets a seller request a currency and lets a fee
  carry a local amount, but never says which amount binds or whether a
  conversion rate is guaranteed. The business must decide what the customer
  pays when the requested currency differs from the airline's pricing
  currency.

  Sources: \href{https://standards.atlassian.net/wiki/spaces/EASD/pages/749305864/Currency+in+Monetary+Amounts}{EASD guide, Currency in Monetary Amounts};
  \href{https://standards.atlassian.net/wiki/spaces/EASD/pages/574685264/Price+Decimals}{EASD guide, Price Decimals};
  \href{https://standards.atlassian.net/wiki/spaces/EASD/pages/574914636/Fees}{EASD guide, Fees};
  \href{https://retailing.iata.org/armi/docs/SHPLOC/}{IATA capability SHPLOC}.
\end{definition}

\begin{definition}[Time limits as UTC instants (C-42)]
  \label{def:offers-time-limits-as-instants}
  \uses{def:instant, thm:offers-offer-expiry-ends-limits, def:offers-payment-limit-at-creation}
  \notready
  Every time limit on an offer or order resolves to one UTC instant before
  it is compared with the current time. Flight times are local to their
  airport and need its offset to resolve.

  Case to check: A payment limit of 2018-01-01 read as midnight UTC and read
  as the end of that day in a seller zone of UTC-12 differ by 36 hours, so
  one payment is on time under one reading and late under the other.

  Open point: The source prefers UTC, but its own time limit examples
  include a bare date and a date and time with no offset. The business must
  decide which instant such a limit denotes, for example the end of that day
  at the point of sale.

  Sources: \href{https://standards.atlassian.net/wiki/spaces/EASD/pages/574488607/Reference+to+Times}{EASD guide, Reference to Times};
  \href{https://standards.atlassian.net/wiki/spaces/EASD/pages/574423066/Time+Limits}{EASD guide, Time Limits}.
\end{definition}
